\documentclass[a4paper,11pt]{article}

% --- PREAMBLE ---
\usepackage[a4paper, top=2.5cm, bottom=2.5cm, left=2cm, right=2cm]{geometry}

\usepackage[T1]{fontenc}
\usepackage[utf8]{inputenc}
\usepackage[italian]{babel}
\usepackage{textcomp}
\usepackage{eurosym}

% Sans Serif
\usepackage{helvet}
\renewcommand{\familydefault}{\sfdefault}

\usepackage{enumitem}
\setlist[itemize]{label=--}

\usepackage{booktabs}
\usepackage{array}
\usepackage{multirow}

% --- PACCHETTI GRAFICA ---
\usepackage{xcolor}
\usepackage{tcolorbox}
\tcbuselibrary{skins, breakable}
\usepackage{hyperref}
\usepackage{titlesec}
\usepackage{amsmath}

% Colori
\definecolor{SmartBlue}{RGB}{0, 70, 140}
\definecolor{SmartGreen}{RGB}{0, 120, 50}
\definecolor{SmartRed}{RGB}{180, 20, 20}
\definecolor{SmartLight}{RGB}{235, 242, 252}

% Hyperref
\hypersetup{
    colorlinks=true,
    linkcolor=SmartBlue,
    urlcolor=SmartBlue,
    pdftitle={Strategia Business e Proiezioni 5 Anni - SmartLight},
    pdfauthor={Team SmartLight Pisa}
}

% Comandi
\newcommand{\pro}[1]{\textcolor{SmartGreen}{\textbf{[+] PRO:} #1}}
\newcommand{\contro}[1]{\textcolor{SmartRed}{\textbf{[-] CONTRO:} #1}}

% Stile sezioni
\titleformat{\section}
  {\color{SmartBlue}\Large\bfseries}
  {\thesection}{1em}{}[\titlerule]

\titleformat{\subsection}
  {\color{SmartBlue}\large\bfseries}
  {\thesubsection}{1em}{}

% --- DOCUMENTO ---
\begin{document}

% --- TITOLO ---
\begin{center}
    \vspace*{1cm}
    {\Huge \bfseries \color{SmartBlue}
        Strategia di Business, Modelli di Revenue \\
        e Proiezioni Finanziarie 2026--2030} \\[0.5cm]
    {\Large \textbf{SmartLight Pisa}} \\[0.2cm]
    {\large \textit{Team SmartLight Pisa}} \\[0.2cm]
    {\small Versione 2.2 --- Marzo 2026}
    \vspace*{0.8cm}
\end{center}

\tableofcontents
\newpage

% ==============================================================
\section*{Obiettivo del Documento}
% ==============================================================
Questo documento analizza i modelli di business per \textbf{SmartLight Pisa},
definisce la struttura dei costi reali per singolo incrocio (CapEx e OpEx),
il pricing SaaS differenziato per tipologia di città, e presenta proiezioni
finanziarie conservative su un orizzonte di 5 anni (2026--2030).

% ==============================================================
\section{Stato Attuale del Progetto (Marzo 2026)}
% ==============================================================

\begin{tcolorbox}[
    enhanced, colback=SmartLight, colframe=SmartBlue,
    title=\textbf{Milestone Completate e In Corso}, boxrule=0.5mm
]
\begin{itemize}[nosep, topsep=4pt]
    \item \textbf{Set 2025 --- Vincita Bando Concorso Idee:} Primo riconoscimento istituzionale. Concept validato.
    \item \textbf{Ott 2025 -- Feb 2026 --- Evoluzione AI:} Sviluppo algoritmi PPO e Dueling Double DQN. Integrazione con SUMO su topologia reale di Pisa. Training $\approx$125\,FPS su GTX 1650.
    \item \textbf{H1 2026 (In corso) --- Fase 3.5 SUMO:} Bridge da simulazione matematica a fisica realistica. Validazione agente su incroci a 4 vie con fasi semaforiche reali, pedoni e veicoli d'emergenza.
    \item \textbf{H2 2026 (Pianificato) --- Pilota Reale:} Deploy su 2 incroci fisici. Dashboard PA. Validazione hardware in campo.
\end{itemize}
\end{tcolorbox}

\vspace{0.3cm}
Il progetto è attualmente in fase \textbf{pre-revenue}. Il focus operativo è la
validazione tecnica necessaria a costruire la prova di valore per avviare
le prime trattative commerciali con la Pubblica Amministrazione.

% ==============================================================
\section{Panoramica dei Modelli di Business}
% ==============================================================

Per una startup DeepTech nel settore Smart City B2G sono stati valutati quattro approcci:

\begin{enumerate}[label=\textbf{\arabic*.}, leftmargin=*, itemsep=1em]

    \item \textbf{Modello Tradizionale (CapEx + Maintenance)}
    \begin{itemize}[label=, leftmargin=1em, nosep]
        \item \textit{Descrizione:} Vendita hardware/software ``chiavi in mano'' + canone manutenzione annuo (10--15\%).
        \item \pro{Incasso immediato (Cash Flow).}
        \item \contro{Ogni vendita richiede gare d'appalto complesse; scarsa ricorrenza; nessun lock-in.}
    \end{itemize}

    \item \textbf{Pure SaaS / TaaS (Traffic as a Service)}
    \begin{itemize}[label=, leftmargin=1em, nosep]
        \item \textit{Descrizione:} Hardware in comodato; il Comune paga abbonamento mensile per intelligenza e dashboard.
        \item \pro{Alta valutazione aziendale (ARR); flusso costante e prevedibile.}
        \item \contro{Difficile contrattualizzazione per PA italiane; elevato fabbisogno di capitale iniziale.}
    \end{itemize}

    \item \textbf{Performance-Based (Risk Sharing)}
    \begin{itemize}[label=, leftmargin=1em, nosep]
        \item \textit{Descrizione:} Divisione dei risparmi generati (carburante, CO$_2$, tempo).
        \item \pro{Argomentazione di vendita fortissima per la PA.}
        \item \contro{Rischio elevatissimo; complessità legale nella certificazione del risparmio.}
    \end{itemize}

    \item \textbf{Licensing / OEM}
    \begin{itemize}[label=, leftmargin=1em, nosep]
        \item \textit{Descrizione:} Vendita dell'algoritmo ``SmartLight Brain'' a grandi player (es.\ Swarco, La Semaforica).
        \item \pro{Scalabilità immediata; zero gestione hardware.}
        \item \contro{Margini ridotti; perdita del brand e del cliente finale.}
    \end{itemize}

\end{enumerate}

% ==============================================================
\section{La Strategia Vincente: ``Hardware-Enabled SaaS''}
% ==============================================================

\subsection*{Approccio ``Land \& Expand'' (Il Cavallo di Troia)}

La strategia adottata è un modello ibrido che abbatte la percezione del rischio
da parte della PA e aggira le complessità burocratiche delle gare d'appalto
nella fase iniziale. Il principio fondamentale è la \textbf{separazione netta}
tra il costo di ingresso (una tantum, Fase 1) e il canone ricorrente (Fase 2):
il primo copre l'hardware, il secondo genera margine.

\begin{tcolorbox}[
    enhanced, colback=blue!5!white, colframe=SmartBlue,
    title=\textbf{La Strategia in 3 Fasi},
    fonttitle=\bfseries\large, boxrule=0.5mm, drop shadow
]

\textbf{\large Fase 1 --- Ingresso ``Shadow Mode'' (Audit \& Monitor)} \\
\textit{Obiettivo: installare l'hardware e validare il valore senza gare d'appalto complesse.}
\begin{itemize}[leftmargin=*, nosep, topsep=5pt]
    \item \textbf{Cosa si vende:} Servizio di audit e monitoraggio del traffico.
    \item \textbf{Prezzo:} Fee una tantum sotto la soglia di affidamento diretto (\texteuro{}5.000--8.000), che \textbf{copre integralmente il CapEx hardware}.
    \item \textbf{Tecnologia:} Kit SmartLight in modalità \textit{shadow}: il sistema non controlla il semaforo, osserva, conta e simula cosa avrebbe fatto l'IA.
    \item \textbf{Value Proposition:} ``Con una spesa minima scoprite quanto CO$_2$ e ore state sprecando ogni giorno --- e quanto risparmiereste.''
\end{itemize}
\vspace{0.3cm}\hrule\vspace{0.3cm}

\textbf{\large Fase 2 --- Attivazione e Controllo (SaaS Ricorrente)} \\
\textit{Obiettivo: convertire il cliente ``passivo'' in un contratto mensile puro margine.}
\begin{itemize}[leftmargin=*, nosep, topsep=5pt]
    \item \textbf{Trigger:} Report Fase 1 (es.\ ``Avete sprecato 400\,h di coda --- l'IA ne avrebbe risparmiate 150'').
    \item \textbf{Cosa si vende:} Controllo attivo (interfaccia API semaforica) + dashboard + manutenzione remota inclusa.
    \item \textbf{Vantaggio:} Nessun nuovo hardware --- si ``accende'' l'intelligenza sull'hardware già installato e già pagato nella Fase 1.
    \item \textbf{Margine:} Il canone SaaS copre solo l'OpEx; il CapEx è già recuperato.
\end{itemize}
\vspace{0.3cm}\hrule\vspace{0.3cm}

\textbf{\large Fase 3 --- Ecosistema di Rete (Data Monetization)} \\
\textit{Obiettivo: creare valore oltre il controllo semaforico.}
\begin{itemize}[leftmargin=*, nosep, topsep=5pt]
    \item \textbf{Soglia:} Con 50+ incroci attivi si possiede la rete di sensori urbani più densa della città.
    \item \textbf{Nuovi flussi:} Dati aggregati e anonimizzati (GDPR) ceduti a logistica, retail, app navigazione, urbanistica predittiva. Esclusivo del Piano B.
    \item \textbf{Edge-first:} Nessun dato identificativo lascia l'infrastruttura locale.
\end{itemize}

\end{tcolorbox}

% ==============================================================
\section{Analisi dei Costi per Incrocio}
% ==============================================================

\subsection{Costo di Setup (CapEx --- Una Tantum per Incrocio)}

\begin{table}[h!]
\centering
\begin{tabular}{>{\bfseries}l l r}
\toprule
Categoria & Dettaglio & Costo \\
\midrule
Hardware AI    & NVIDIA Jetson Orin (edge computing)           & \texteuro{}800 \\
Visione        & 2$\times$ Telecamere IP 4K H.265               & \texteuro{}600 \\
Connettività   & Router LTE industriale                         & \texteuro{}400 \\
Protezione     & Box stagno IP65 + cablaggio + fissaggi         & \texteuro{}300 \\
Installazione  & Squadra tecnica + noleggio cestello elevatore  & \texteuro{}1.300 \\
Setup AI       & Fine-tuning modello, mapping ROI, validazione  & \texteuro{}1.300 \\
\midrule
\multicolumn{2}{l}{\textbf{TOTALE CapEx per Incrocio}} & \textbf{\texteuro{}4.700} \\
\bottomrule
\end{tabular}
\caption{Costo di installazione per singolo incrocio medio. Coperto dalla fee una tantum della Fase 1.}
\end{table}

\subsection{Costi Operativi Ricorrenti (OpEx --- Per Incrocio / Anno)}

L'hardware SmartLight è \textbf{solid-state} (Jetson, camere IP): nessuna parte meccanica soggetta a usura rapida. Il modello AI viene aggiornato periodicamente e distribuito via push remoto sull'intera flotta con costo fisso di R\&D (già nel budget team), trascurabile come quota per singolo incrocio. La connettività è una SIM IoT a bassa banda condivisa su piano aziendale multi-nodo, non un abbonamento individuale.

\begin{table}[h!]
\centering
\begin{tabular}{>{\bfseries}l l r}
\toprule
Categoria & Dettaglio & Costo/Anno \\
\midrule
Manutenzione hardware & Visita preventiva ogni $\approx$2 anni + ricambi         & \texteuro{}150 \\
Connettività          & SIM IoT low-bandwidth, quota piano condiviso multi-nodo  & \texteuro{}70  \\
Cloud / monitoring    & Storage metriche, monitoring remoto, backup              & \texteuro{}60  \\
Aggiornamenti SW/AI   & Quota deploy remoto (R\&D a costo fisso aziendale)       & \texteuro{}50  \\
Supporto tecnico      & Quota helpdesk remoto (amortizzato su flotta)            & \texteuro{}40  \\
\midrule
\multicolumn{2}{l}{\textbf{TOTALE OpEx per Incrocio / Anno}} & \textbf{\texteuro{}370} \\
\bottomrule
\end{tabular}
\caption{Costi operativi annui per incrocio a regime. I costi fissi aziendali (R\&D, sales, team) sono separati.}
\end{table}

% ==============================================================
\section{Pricing SaaS verso le Città}
% ==============================================================

SmartLight adotta \textbf{due soli piani commerciali} per semplicità negoziale
e chiarezza nella proposta alla PA. Il principio è che il canone SaaS deve
essere \textbf{immediatamente sostenibile} per qualsiasi Comune, rendendo
il ROI evidente anche ai più scettici. Il CapEx è già coperto dalla Fase 1.

\begin{tcolorbox}[
    enhanced, colback=blue!5!white, colframe=SmartBlue,
    title=\textbf{I Due Piani SmartLight}, fonttitle=\bfseries\large,
    boxrule=0.5mm, drop shadow
]

\textbf{\large Piano A --- Città Piccola}
{\small (Comuni con meno di 50.000 abitanti, tipicamente 3--10 incroci)}

\begin{tabular}{>{\bfseries}l p{9cm}}
Canone              & \texteuro{}100 / incrocio / mese \\
Include             & AI controllo semaforico adattivo, dashboard base (report mensile PDF), manutenzione remota inclusa \\
Escluso             & Priorità bus/ambulanze, API dati, data monetization \\
Fee Fase 1 (CapEx)  & Pacchetto fisso \texteuro{}5.000--7.000 (copre installazione hardware) \\
Esempio (5 incroci) & \texteuro{}6.000 / anno di canone SaaS \\
\end{tabular}

\vspace{0.4cm}\hrule\vspace{0.4cm}

\textbf{\large Piano B --- Città Grande}
{\small (Comuni con più di 50.000 abitanti, tipicamente 20+ incroci)}

\begin{tabular}{>{\bfseries}l p{9cm}}
Canone              & \texteuro{}150 / incrocio / mese \\
Include             & AI controllo adattivo, priorità Bus/Ambulanze, dashboard real-time, API dati, SLA 99.5\%, manutenzione inclusa, data monetization (dal 2028) \\
Fee Fase 1 (CapEx)  & \texteuro{}4.700 / incrocio (scalato in proporzione al volume) \\
Esempio (25 incroci) & \texteuro{}45.000 / anno di canone SaaS \\
\end{tabular}

\end{tcolorbox}

\vspace{0.3cm}
\noindent\textbf{Razionale del pricing:} prezzi volutamente accessibili per abbattere
la barriera di ingresso nella PA italiana. Il margine nasce dal \textbf{volume}
(molti incroci per città, molte città) e dalla struttura a costi fissi bassi
(hardware già pagato in Fase 1, OpEx €370/anno).

\subsection{ROI per il Comune --- Argomento di Vendita}

\begin{tcolorbox}[
    enhanced, colback=green!4!white, colframe=SmartGreen,
    title=\textbf{Piano A --- 5 Incroci, Città Piccola}, boxrule=0.5mm
]
\begin{tabular}{lr}
Costo SaaS annuo ($100 \times 12 \times 5$) & \textbf{\texteuro{}6.000 / anno} \\
\midrule
Risparmio carburante collettivo ($-25\%$ attese) & $\approx$\texteuro{}30.000 / anno \\
CO$_2$ ridotta ($5 \times 87\,\text{t}$) & $\approx$435\,t CO$_2$ / anno \\
\midrule
\textbf{ROI netto per il Comune} & $\approx$ \textbf{5$\times$ il costo SaaS} \\
\end{tabular}
\end{tcolorbox}

\begin{tcolorbox}[
    enhanced, colback=green!4!white, colframe=SmartGreen,
    title=\textbf{Piano B --- 25 Incroci, Città Grande}, boxrule=0.5mm
]
\begin{tabular}{lr}
Costo SaaS annuo ($150 \times 12 \times 25$) & \textbf{\texteuro{}45.000 / anno} \\
\midrule
Risparmio carburante collettivo ($-25\%$ attese) & $\approx$\texteuro{}150.000 / anno \\
CO$_2$ ridotta ($25 \times 87\,\text{t}$) & $\approx$2.175\,t CO$_2$ / anno \\
Valore ESG / accesso fondi Green Deal EU & non quantificato \\
\midrule
\textbf{ROI netto per il Comune} & $\approx$ \textbf{3.3$\times$ il costo SaaS} \\
\end{tabular}
\end{tcolorbox}

\noindent Il canone SaaS non è una spesa --- è un investimento con ritorno misurabile
e certificabile. Questo è il cardine dell'argomentazione nelle trattative con
gli assessori alla mobilità e gli uffici acquisti pubblici.

\subsection{Margine Netto per Incrocio a Regime}

Poiché il CapEx è recuperato interamente tramite la fee una tantum della Fase 1,
il canone SaaS è \textbf{quasi interamente margine}, al netto del solo OpEx annuo.

\begin{table}[h!]
\centering
\renewcommand{\arraystretch}{1.3}
\begin{tabular}{l r r}
\toprule
\textbf{Voce} & \textbf{Piano A (\texteuro{}100/mese)} & \textbf{Piano B (\texteuro{}150/mese)} \\
\midrule
Ricavo SaaS annuo / incrocio          & \texteuro{}1.200 & \texteuro{}1.800 \\
OpEx annuo                            & $-$\texteuro{}370 & $-$\texteuro{}370 \\
CapEx (coperto da Fase 1)             & \texteuro{}0     & \texteuro{}0     \\
\midrule
\textbf{Margine netto / incrocio / anno} & \textbf{\texteuro{}830 (69\%)} & \textbf{\texteuro{}1.430 (79\%)} \\
\bottomrule
\end{tabular}
\caption{Economics per singolo incrocio a regime. Il CapEx è già recuperato tramite la fee della Fase 1, quindi non pesa sul canone ricorrente.}
\end{table}

% ==============================================================
\section{Proiezioni Finanziarie 2026--2030}
% ==============================================================

\subsection{Assunzioni del Modello}

\begin{itemize}[nosep, topsep=4pt]
    \item Il pilota reale (2 incroci) parte H2 2026 in shadow mode. Nessun ricavo SaaS nel 2026.
    \item Dal 2027 inizia la vendita SaaS piena. Mix stimato: 50\% Piano A / 50\% Piano B, prezzo medio ponderato \textbf{\texteuro{}125/mese/incrocio}.
    \item La \textbf{data monetization} (solo Piano B) è attivata a partire dal 2028, al raggiungimento di 50+ incroci totali.
    \item I \textbf{ricavi una tantum} (Fase 1, audit + CapEx recovery) sono separati dal SaaS e calcolati conservativamente.
    \item Il CapEx nuovi incroci è coperto dai ricavi Fase 1 o da bandi PNRR; non pesa sul FCF operativo.
    \item I costi R\&D/Team crescono in modo controllato: team small (2--4 persone nel 2026--2027), poi graduale scalata.
\end{itemize}

\subsection{Tabella di Proiezione (Scenario Base --- Conservativo)}

\begin{table}[h!]
\centering
\renewcommand{\arraystretch}{1.5}
\begin{tabular}{l r r r r r}
\toprule
\textbf{Voce} & \textbf{2026} & \textbf{2027} & \textbf{2028} & \textbf{2029} & \textbf{2030} \\
\midrule
\textbf{Incroci Attivi (SaaS)}        & 0          & 15         & 50          & 150         & 350         \\
\midrule
Ricavo SaaS (\texteuro{}125/mese avg.)  & \texteuro{}0 & \texteuro{}23k & \texteuro{}75k & \texteuro{}225k & \texteuro{}525k \\
Ricavo Data Monetization (Piano B)    & \texteuro{}0 & \texteuro{}0   & \texteuro{}10k & \texteuro{}45k  & \texteuro{}105k \\
Ricavo Fase 1 (una tantum, CapEx)     & \texteuro{}0 & \texteuro{}30k & \texteuro{}50k & \texteuro{}90k  & \texteuro{}130k \\
\midrule
\textbf{Ricavi Totali}                & \texteuro{}0   & \texteuro{}53k & \texteuro{}135k & \texteuro{}360k & \texteuro{}760k \\
\midrule
OpEx incroci (\texteuro{}370/incr.)   & \texteuro{}3k  & \texteuro{}6k  & \texteuro{}18k  & \texteuro{}55k  & \texteuro{}130k \\
Costi R\&D \& Team                    & \texteuro{}22k & \texteuro{}50k & \texteuro{}75k  & \texteuro{}110k & \texteuro{}180k \\
Marketing \& Business Dev.            & \texteuro{}0   & \texteuro{}5k  & \texteuro{}12k  & \texteuro{}25k  & \texteuro{}45k  \\
\midrule
\textbf{Costi Totali}                 & \texteuro{}25k & \texteuro{}61k & \texteuro{}105k & \texteuro{}190k & \texteuro{}355k \\
\midrule
\textbf{Free Cash Flow}               &
    \textcolor{SmartRed}{\textbf{$-$\texteuro{}25k}} &
    \textcolor{SmartRed}{\textbf{$-$\texteuro{}8k}} &
    \textcolor{SmartGreen}{\textbf{$+$\texteuro{}30k}} &
    \textcolor{SmartGreen}{\textbf{$+$\texteuro{}170k}} &
    \textcolor{SmartGreen}{\textbf{$+$\texteuro{}405k}} \\
\bottomrule
\end{tabular}
\caption{Proiezioni finanziarie 2026--2030 (scenario base, arrotondamenti al migliaio). I ricavi Fase 1 compensano il CapEx nuovi incroci. ARR 2030: \texteuro{}630k (SaaS + Data). FCF cumulato 2026--2030: $\approx$\texteuro{}570k.}
\end{table}

\begin{tcolorbox}[
    enhanced, colback=green!5!white, colframe=SmartGreen,
    title=\textbf{Milestone Finanziarie Chiave}, boxrule=0.5mm
]
\begin{itemize}[nosep, topsep=4pt]
    \item \textbf{2026:} Pre-revenue. Costo netto $\approx$\texteuro{}25k (team + pilota + preparazione commerciale).
    \item \textbf{2027:} Primi ricavi reali (\texteuro{}53k). Disavanzo SaaS quasi nullo ($-$\texteuro{}8k) grazie ai ricavi Fase 1.
    \item \textbf{2028:} \textbf{Breakeven operativo.} $+$\texteuro{}30k FCF. 50 incroci = 3--5 città attive.
    \item \textbf{2029:} Crescita accelerata. 150 incroci $\approx$ 8--10 città. FCF $+$\texteuro{}170k.
    \item \textbf{2030:} ARR \texteuro{}630k. FCF $+$\texteuro{}405k. Modello pienamente scalato e autofinanziato.
\end{itemize}
\end{tcolorbox}

\subsection{Sensitivity Analysis --- Tre Scenari al 2030}

\begin{table}[h!]
\centering
\renewcommand{\arraystretch}{1.3}
\begin{tabular}{l r r r}
\toprule
\textbf{Scenario} & \textbf{Incroci 2030} & \textbf{ARR 2030} & \textbf{FCF 2030} \\
\midrule
\textbf{Pessimistico} (burocrazia lenta, 1--2 città/anno) & 100  & \texteuro{}225k & \texteuro{}60k  \\
\textbf{Base} (3--5 città/anno, modello attuale)           & 350  & \texteuro{}630k & \texteuro{}405k \\
\textbf{Ottimistico} (PNRR accelerato, 6--10 città/anno)  & 700  & \texteuro{}1.26M & \texteuro{}900k \\
\midrule
\multicolumn{4}{l}{\small Leva principale: velocità di approvazione PA e capacità del team Business Development.} \\
\bottomrule
\end{tabular}
\caption{Sensitivity analysis su tre scenari di crescita al 2030.}
\end{table}

\noindent\textit{Nota sullo scenario pessimistico:} anche con soli 100 incroci attivi al 2030, il modello rimane profittevole (\texteuro{}60k FCF) e autosufficiente. Questo è il pavimento di sicurezza del business.

% ==============================================================
\section{Analisi della Sostenibilità Legale (Italia)}
% ==============================================================

\begin{description}[style=multiline, leftmargin=4.5cm, font=\bfseries\color{SmartBlue}]
    \item[Affidamento Diretto] La Fase 1 (Audit) rientra ampiamente sotto la soglia dei \texteuro{}140.000 (art.\ 50 D.Lgs.\ 36/2023), riducendo i tempi di acquisizione da anni a settimane. Anche il canone SaaS annuo di una città piccola (es.\ \texteuro{}6.000) rimane abbondantemente sotto soglia.
    \item[Spesa Corrente IT] Il canone SaaS mensile è classificabile come ``spesa corrente per servizi IT'', categoria già presente nei bilanci comunali, evitando l'iter (più complesso) dell'investimento infrastrutturale pluriennale.
    \item[PNRR / Green Deal] La compliance ESG (CO$_2$ certificabile, mobilità sostenibile M2C2) allinea la spesa a KPI obbligatori per l'accesso ai fondi europei, rendendo SmartLight un facilitatore di finanziamenti.
    \item[GDPR Compliance] Il modello edge-first garantisce che nessun dato personale lasci l'infrastruttura locale, semplificando l'approvazione del DPO comunale e riducendo il rischio di contestazioni privacy.
    \item[Project Financing] Per estensioni su larga scala ($>$50 incroci), è proponibile un Project Financing: investimento privato iniziale, Comune paga canone su performance misurate.
\end{description}

% ==============================================================
\section{Conclusioni e Next Steps}
% ==============================================================

SmartLight presenta un modello economico solido e difendibile: il CapEx è coperto
dalla fee di ingresso in Fase 1, il canone SaaS è accessibile per qualsiasi PA
(\texteuro{}100--150/mese/incrocio), il ROI per il Comune è 3--5$\times$ dimostrabile
con dati reali, e anche nello scenario pessimistico il modello resta profittevole.
La traiettoria verso i \texteuro{}600k+ di ARR al 2030 non dipende da megacontratti,
ma dall'accumulo progressivo di contratti medio-piccoli tramite affidamento diretto.

\vspace{0.4cm}
\noindent\textbf{Azioni prioritarie per il 2026:}
\begin{enumerate}[nosep, topsep=4pt]
    \item Completare il pilota reale su 2 incroci (H2 2026) e pubblicare i risultati misurati.
    \item Rilasciare la dashboard PA come strumento mostrabile nelle trattative commerciali.
    \item Identificare 3--5 Comuni target per contratti Fase 1 (Audit) entro fine 2026.
    \item Esplorare bandi PNRR ``Smart Cities'' per co-finanziare il CapEx dei primi 10--15 incroci.
    \item Strutturare il primo accordo pilota di data monetization (anche informale).
\end{enumerate}

\vspace{1cm}
\begin{center}
    \textit{``SmartLight non chiede alle città di credere nella tecnologia ---}\\
    \textit{chiede di leggere i dati del proprio incrocio.''}
\end{center}

\end{document}